% $HeadURL$

\subsection{Glyph: \glyph{Perturbing agent}}
\label{sec:perturbing agent}
 
Biochemical networks can be affected by external influences.  Those
influences can be the effect of well-defined physical perturbing agents, such as a light
pulse or a change in temperature; they can also be more complex and not
well-defined phenomena, for instance the outcome of a biological process, an experimental
setup, or a mutation.  For these situations, SBGN provides the
\glyph{perturbing agent} glyph. It is an EPN, and represents the amount to perturbing agent applied to a process.  

\begin{glyphDescription}

\glyphSboTerm SBO:0000405 ! perturbing agent

\glyphContainer A \glyph{perturbing agent} is represented by a modified hexagon
having two opposite concave faces, as illustrated in \fig{perturbing agent}.

\glyphLabel A \glyph{perturbing agent} is identified by a label placed in an
unbordered box containing a string of characters.  The characters can be
distributed on several lines to improve readability, although this is not
mandatory.  The label box must be attached to the center of the
\glyph{perturbing agent} container.  The label may spill outside of the container.

\glyphAux A \glyph{perturbing agent} may carry a \glyph{clone marker}
(\sect{cloneMarker}).

A \glyph{perturbing agent} can optionally carry one or more \glyph{units of information} (\sect{unitInfo}).  A controlled vocaulary is available to describing the physical characteristic of the perturbing agent (see \sect{physical-characteristics-cv}).  

\end{glyphDescription}

\begin{figure}[H]
  \centering
  \includegraphics[scale = 0.3]{images/perturbing_agent}
  \caption{The \PD glyph for \glyph{perturbing agent}.}
  \label{fig:perturbing agent}
\end{figure}




% The following is for [X]Emacs users.  Please leave in place.
% Local Variables:
% TeX-master: "../sbgn_PD-level1"
% End:

